\documentclass[10pt,a4paper]{amsart}
\usepackage{amsmath,amssymb,amsthm,graphicx}
\usepackage{verbatim}
\newenvironment{reference}{\comment}{\endcomment}
%\newenvironment{reference}{}{}
\newenvironment{slogan}{\comment}{\endcomment}
\newenvironment{history}{\comment}{\endcomment}

% For commutative diagrams we use Xy-pic
\usepackage[all]{xy}

% We use 2cell for 2-commutative diagrams.
\xyoption{2cell}
% Theorem environments.
%
\theoremstyle{plain}
\newtheorem{theorem}[subsection]{Theorem}
\newtheorem{proposition}[subsection]{Proposition}
\newtheorem{lemma}[subsection]{Lemma}

\theoremstyle{definition}
\newtheorem{definition}[subsection]{Definition}
\newtheorem{example}[subsection]{Example}
\newtheorem{exercise}[subsection]{Exercise}
\newtheorem{situation}[subsection]{Situation}

\theoremstyle{remark}
\newtheorem{remark}[subsection]{Remark}
\newtheorem{remarks}[subsection]{Remarks}

\numberwithin{equation}{subsection}

% Macros
%
\def\lim{\mathop{\mathrm{lim}}\nolimits}
\def\colim{\mathop{\mathrm{colim}}\nolimits}
\def\Spec{\mathop{\mathrm{Spec}}}
\def\Hom{\mathop{\mathrm{Hom}}\nolimits}
\def\Ext{\mathop{\mathrm{Ext}}\nolimits}
\def\SheafHom{\mathop{\mathcal{H}\!\mathit{om}}\nolimits}
\def\SheafExt{\mathop{\mathcal{E}\!\mathit{xt}}\nolimits}
\def\Sch{\mathit{Sch}}
\def\Mor{\mathop{\mathrm{Mor}}\nolimits}
\def\Ob{\mathop{\mathrm{Ob}}\nolimits}
\def\Sh{\mathop{\mathit{Sh}}\nolimits}
\def\NL{\mathop{N\!L}\nolimits}
\def\CH{\mathop{\mathrm{CH}}\nolimits}
\def\proetale{{pro\text{-}\acute{e}tale}}
\def\etale{{\acute{e}tale}}
\def\QCoh{\mathit{QCoh}}
\def\Ker{\mathop{\mathrm{Ker}}}
\def\Im{\mathop{\mathrm{Im}}}
\def\Coker{\mathop{\mathrm{Coker}}}
\def\Coim{\mathop{\mathrm{Coim}}}

% Boxtimes
%
\DeclareMathSymbol{\boxtimes}{\mathbin}{AMSa}{"02}

%
% Macros for moduli stacks/spaces
%
\def\QCohstack{\mathcal{QC}\!\mathit{oh}}
\def\Cohstack{\mathcal{C}\!\mathit{oh}}
\def\Spacesstack{\mathcal{S}\!\mathit{paces}}
\def\Quotfunctor{\mathrm{Quot}}
\def\Hilbfunctor{\mathrm{Hilb}}
\def\Curvesstack{\mathcal{C}\!\mathit{urves}}
\def\Polarizedstack{\mathcal{P}\!\mathit{olarized}}
\def\Complexesstack{\mathcal{C}\!\mathit{omplexes}}
% \Pic is the operator that assigns to X its picard group, usage \Pic(X)
% \Picardstack_{X/B} denotes the Picard stack of X over B
% \Picardfunctor_{X/B} denotes the Picard functor of X over B
\def\Pic{\mathop{\mathrm{Pic}}\nolimits}
\def\Picardstack{\mathcal{P}\!\mathit{ic}}
\def\Picardfunctor{\mathrm{Pic}}
\def\Deformationcategory{\mathcal{D}\!\mathit{ef}}
\usepackage{cleveref}
\crefname{originallemma}{original lemma}{original lemmas}
\Crefname{originallemma}{Original lemma}{Original lemmas}
\makeatletter
\newlabel{algebra-lemma-characterize-integral-element}{{052I}{source}{}{sourcepointer-1}{}}
\newlabel{algebra-lemma-characterize-integral-element@cref}{{[originallemma][0][]052I}{[1][0][]source}{}{}{}}
\makeatother
\title{A complete ring with a principal ideal whose completed localization is not flat}
\author{The Stacks Project Authors}
\date{stacks-project:examples.tex:local-snapshot-20261012; GNU FDL 1.2 or later}
\begin{document}
\maketitle
\noindent\textit{Editorial note.} All delivered mathematical text is the author's own, copied from the pinned Stacks Project chapter file examples.tex; no statement or proof step was rewritten, reordered or removed.\par\medskip
\noindent\textit{Editorial note.} The counted proof is the author's complete argument: The counted proof is the two-line instantiation; the load-bearing non-flatness lives complete in the context span (kernel comparison + integrality absurdity); 'observe' bridges via R[[x]]\_f = R\_f[[x]] and x-adic completeness of power series rings.\par\medskip
\noindent\textit{Editorial note.} The label lemma-nonflat-completion-localization names the counted statement, and the proof environment immediately following it in the source is delivered verbatim as its proof.\par\medskip
\noindent\textit{Editorial note.} The counted statement is carried by the author's own lemma environment in the chapter source.\par\medskip
\noindent\textit{Allowed standard prerequisites (editorial).}\begin{itemize}
\item Context span 1103-1140: with R = k[y, z, a\_1, a\_2, ...]/(ya\_i, a\_i a\_j) and f the residue class of z, the ring map R[[x]] -> R\_f[[x]] is not flat - the kernels I of y : R[[x]] -> R[[x]] and J of y : R\_f[[x]] -> R\_f[[x]] satisfy J != I R\_f[[x]], proved by comparing the coefficient of a\_n after reduction mod (y, a\_1, ...), Noetherian finiteness in k[z][[x]], and Algebra lemma-characterize-integral-element giving the absurd integrality of z\textasciicircum{}\{-1\}x; flatness would force J = I R\_f[[x]] by tensoring 0 -> I -> R[[x]] -> R[[x]] -> 0
\item x-adic completion of power series rings: R[[x]] is x-adically complete, and localizing at a coefficient-ring element commutes with power series: R[[x]]\_f = R\_f[[x]]
\end{itemize}\medskip
\ifdefined\hypertarget\hypertarget{sourcepointer-1}{}\fi
\noindent\textit{Source reference.} Original lemma 052I: sources/stacks-chow/tags.txt, source lines 6571--6571. Prerequisite pointer invoked at 1183; it names a standard fact of a sibling chapter that the counted argument uses, and every step of the counted proof built on it is retained verbatim in the delivered body.\par\medskip
Next, we will construct an example where the completion of a localization
is nonflat. To do this consider the ring
$$
R = k[y, z, a_1, a_2, a_3, \ldots]/(ya_i, a_i a_j)
$$
Denote $f \in R$ the residue class of $z$. We claim the ring map
\begin{equation}
\label{equation-nonflat}
R[[x]] \longrightarrow R_f[[x]]
\end{equation}
isn't flat. Let $I$ be the kernel of $y : R[[x]] \to R[[x]]$. A typical
element $g$ of $I$ looks like $g = \sum g_{n, m} a_mx^n$
where $g_{n, m} \in k[z]$ and for a given $n$ only a finite number of
nonzero $g_{n, m}$. Let $J$ be the kernel of $y : R_f[[x]] \to R_f[[x]]$.
We claim that $J \not = I R_f[[x]]$. Namely, if this were true then we
would have
$$
\sum z^{-n} a_n x^n = \sum\nolimits_{i = 1, \ldots, m} h_i g_i
$$
for some $m \geq 1$, $g_i \in I$, and $h_i \in R_f[[x]]$. Say
$h_i = \bar h_i \bmod (y, a_1, a_2, a_3, \ldots)$
with $\bar h_i \in k[z, 1/z][[x]]$.  Looking at the coefficient of
$a_n$ and using the description of the elements $g_i$ above we would get
$$
z^{-n} x^n = \sum \bar h_i \bar g_{i, n}
$$
for some $\bar g_{i, n} \in k[z][[x]]$. This would mean that
all $z^{-n}x^n$ are contained in the finite $k[z][[x]]$-module
generated by the elements $\bar h_i$. Since $k[z][[x]]$ is Noetherian
this implies that the $R[z][[x]]$-submodule of $k[z, 1/z][[x]]$
generated by $1, z^{-1}x, z^{-2}x^2, \ldots$ is finite. By
Algebra, Lemma \ref{algebra-lemma-characterize-integral-element}
we would conclude that $z^{-1}x$ is integral over $k[z][[x]]$
which is absurd. On the other hand,
if (\ref{equation-nonflat}) were flat, then we would
get $J = IR_f[[x]]$ by tensoring the exact sequence
$0 \to I \to R[[x]] \xrightarrow{y} R[[x]]$ with $R_f[[x]]$.


\begin{lemma}
\label{lemma-nonflat-completion-localization}
There exists a ring $A$ complete with respect to a principal ideal $I$
and an element $f \in A$ such that the $I$-adic completion
$A_f^\wedge$ of $A_f$ is not flat over $A$.
\end{lemma}

\begin{proof}
Set $A = R[[x]]$ and $I = (x)$ and observe that $R_f[[x]]$
is the completion of $R[[x]]_f$.
\end{proof}

\end{document}
